% Generated draft: requires mathematical and source review.
\documentclass[11pt]{article}
\usepackage[margin=25mm]{geometry}
\usepackage{fontspec}
\setmainfont{Latin Modern Roman}
\usepackage{amsmath,amssymb,mathtools,unicode-math}
\setmathfont{Latin Modern Math}
\usepackage{enumitem,needspace,xurl,hyperref}
\setlength{\parindent}{0pt}
\setlength{\parskip}{5pt}
\setlength{\emergencystretch}{3em}
\newcommand{\N}{\mathbb{N}}
\newcommand{\Z}{\mathbb{Z}}
\newcommand{\Q}{\mathbb{Q}}
\newcommand{\R}{\mathbb{R}}
\newcommand{\C}{\mathbb{C}}
\newcommand{\F}{\mathbb{F}}
\newcommand{\A}{\mathbb{A}}
\newcommand{\PP}{\mathbb{P}}
\newcommand{\E}{\mathbb{E}}
\newcommand{\eps}{\varepsilon}
\newcommand{\dd}{\,\mathrm{d}}
\begin{document}
\section*{UnsolvedMath ID: 700298}
\textbf{Source problem code:} catalogue entry
\subsection{Definitions and mathematical statement}
Let \(G\) be a finite group, and let \(\operatorname{Syl}_2(G)\) denote the set of its Sylow \(2\)-subgroups. For \(X\leq G\), define its normalizer by
\[
N_G(X)=\{g\in G:gXg^{-1}=X\}.
\]
A \(2\)-subgroup is a subgroup of order a power of \(2\). We use the standard convention that a \(2\)-local subgroup of \(G\) is \(N_G(U)\) for a nontrivial \(2\)-subgroup \(U\leq G\); the supplied record does not explicitly define this term. Allowing \(U=1\) would not change the condition below, since \([G:N_G(1)]=1\).

A \emph{Sylow intersection} is a subgroup
\[
I=\bigcap_{P\in\mathcal S}P,
\qquad \varnothing\neq\mathcal S\subseteq\operatorname{Syl}_2(G).
\]
A \emph{paired Sylow intersection} is \(P\cap Q\), where \(P,Q\in\operatorname{Syl}_2(G)\). The source does not specify whether the family may be empty or whether a pair must have distinct members. We use nonempty families and allow \(P=Q\). These choices do not affect the requested classifications: an empty intersection, interpreted as \(G\), has normalizer of index \(1\), while the normalizer of a Sylow \(2\)-subgroup always has odd index.

For each of the following four conditions separately, describe all finite groups \(G\) satisfying that condition:
\begin{enumerate}
\item[(a)] Every \(2\)-local subgroup has odd index in \(G\); equivalently,
\[
\forall U\leq G,\qquad
\bigl(|U|=2^k\text{ for some integer }k\geq1\bigr)
\Longrightarrow [G:N_G(U)]\text{ is odd}.
\]
\item[(b)] Every normalizer of a Sylow intersection has odd index in \(G\):
\[
\forall\varnothing\neq\mathcal S\subseteq\operatorname{Syl}_2(G),\qquad
\left[G:N_G\!\left(\bigcap_{P\in\mathcal S}P\right)\right]\text{ is odd}.
\]
\item[(c)] Every normalizer of a paired Sylow intersection has odd index in \(G\):
\[
\forall P,Q\in\operatorname{Syl}_2(G),\qquad
[G:N_G(P\cap Q)]\text{ is odd}.
\]
\item[(d)] For every \(P,Q\in\operatorname{Syl}_2(G)\), writing \(\langle P,Q\rangle\) for the subgroup generated by \(P\) and \(Q\), there exists
\[
S\in\operatorname{Syl}_2(\langle P,Q\rangle)
\quad\text{such that}\quad P\cap Q\trianglelefteq S.
\]
Here normality includes the requirement \(P\cap Q\leq S\).
\end{enumerate}
No assumption that \(|G|\) is even is imposed. In particular, the conditions include groups of odd order. The supplied registry labels the problem open; this is a supplied status, not an independently verified mathematical or literature assessment.
\subsection{Short English statement}
Classify, separately, the finite groups satisfying each of four conditions involving Sylow \(2\)-subgroups: odd index for all \(2\)-local subgroups, odd index for normalizers of all Sylow intersections, odd index for normalizers of intersections of two Sylow subgroups, and normality of each such paired intersection in some Sylow \(2\)-subgroup of the subgroup generated by the pair.
\subsection{Sources}
\begin{itemize}
\item Local UnsolvedMath registry, supplied complete source record: permanent ID \(700298\), exact source problem code \texttt{catalogue entry}, title \emph{Kourovka Notebook Problem 5.14}. The record supplies the status \texttt{open} and an update timestamp of \texttt{2026-10-16T14:00:00Z}; these are registry data, not independent verification.
\item \emph{The Kourovka Notebook (Unsolved Problems in Group Theory)}, original Problem 5.14, 5th issue (1976); the supplied background identifies the 21st issue edition (2026), \texttt{arXiv:1401.0300v48}. The collection description names E. I. Khukhro and V. D. Mazurov as editors.
\item Proposed by V. V. Kabanov, A. A. Makhn\"ev, and A. I. Starostin; proposed year 1976, as supplied in the record.
\end{itemize}
\end{document}
